\documentclass[11pt,a4paper]{article}
\usepackage[T1]{fontenc}
\usepackage{isabelle,isabellesym}
\usepackage{amssymb}
\usepackage{stmaryrd}
\usepackage{wasysym}
\usepackage{textcomp}
% \isasymbool renders as \bbbB, which needs the txmia font from txfonts
% (isabellesym.sty:154). Loaded after amssymb so its symbol definitions win.
\usepackage{txfonts}

% `\<^enum>` is a DOCUMENT-group symbol (etc/symbols: code 0x25b8), meant for
% list markup inside text comments, where Isabelle's markup handler renders it
% and no macro is needed. ModelKit/Tagging.thy repurposes it as the mixfix of
% the `Part` operator, i.e. inside TERMS -- and there the symbol goes out as
% `\isactrlenum`, which nothing defines: not in Isabelle's own texinputs (the
% only place isabelle.sty lives; the system TeX Live has no isabelle*.sty at
% all), not anywhere in the Isabelle2026-RC0 tree. Define it here with the
% glyph the symbol table assigns, so the printed operator matches what jEdit
% shows.
\newcommand{\isactrlenum}{\ensuremath{\blacktriangleright}}

% this should be the last package used
\usepackage{pdfsetup}

% urls in roman style, theory text in math-similar italics
\urlstyle{rm}
\isabellestyle{it}

\begin{document}

\title{Isabelle/HOL/GST: A Formal Proof Environment\\ for Generalized Set Theories}
\author{Ciar\'an Dunne \and J. B. Wells}

\maketitle

\begin{abstract}
  A \emph{generalized set theory} (GST) is like a standard set theory but
  also has non-set structured objects that can contain other structured
  objects including sets.  This development provides Isabelle/HOL support for
  GSTs, which are treated as type classes that combine \emph{features}
  specifying kinds of mathematical objects, e.g., sets, ordinal numbers,
  ordered pairs, functions, natural numbers, and an exception object.

  Specialized type-like predicates called \emph{soft types} are used
  throughout.  When a GST is assembled from features, extra axioms are
  generated following a user-modifiable policy to fill specification gaps:
  that every object is owned by exactly one feature, which combinations of
  objects of different features are allowed, and what each operation returns
  outside its intended domain.

  Although a GST can be used without a model, for confidence in its
  consistency a model is built for each GST from components that specify each
  feature's contribution to each tier of a von-Neumann-style cumulative
  hierarchy defined via ordinal recursion; that model is then connected to a
  separate type which the GST occupies, using Isabelle's Lifting and Transfer
  packages.

  The development bootstraps from Paulson's \emph{ZFC in HOL}, instantiating
  the set, ordinal, ordinal-recursion, ordered-pair and function features at
  its type \isa{V}, which yields a GST in which every object is a set.  It
  then assembles the example GSTs of the two papers below and builds a model
  for one of them at a separate type.
\end{abstract}

\tableofcontents

\parindent 0pt\parskip 0.5ex

\section{Sources}

The development accompanies:

\begin{itemize}
  \item Ciar\'an Dunne and J. B. Wells, \emph{Isabelle/HOL/GST: A Formal
    Proof Environment for Generalized Set Theories}, 2022.
    \url{https://arxiv.org/abs/2207.12039}
  \item Ciar\'an Dunne, J. B. Wells and Fairouz Kamareddine,
    \emph{Generating Custom Set Theories with Non-Set Structured Objects},
    CICM 2021, LNCS 12833, pp. 228--244.
    \url{https://doi.org/10.1007/978-3-030-81097-9_19}
\end{itemize}

\input{session}

\end{document}
